% =====================================================================
% Structured Task Decomposition Improves Reliability of LLM-Based
% Knowledge Analysis
%
% Standard article class, arXiv-friendly. Compile with pdflatex +
% bibtex + pdflatex x2.
%
% NOTE FOR THE AUTHOR (remove before submission):
%   - Replace the \author{...} placeholder with real name/affiliation/email.
%   - This file is the LaTeX rendering of paper/reliability-benchmark.md.
%   - All experimental numbers (31%/92%/+8/+54, n=13/n=7, 7/7, 3/3,
%     3.14/4.14, etc.) are transcribed verbatim from the source report;
%     do not round or alter them.
% =====================================================================
\documentclass[11pt]{article}

% ----- page geometry -----
\usepackage[margin=1in]{geometry}

% ----- typography -----
\usepackage[T1]{fontenc}
\usepackage[utf8]{inputenc}
\usepackage{microtype}
\usepackage{lmodern}

% ----- math -----
\usepackage{amsmath}
\usepackage{amssymb}

% ----- tables -----
\usepackage{booktabs}
\usepackage{array}
\usepackage{tabularx}

% ----- code / verbatim -----
\usepackage{verbatim}
\usepackage{listings}
\lstset{
  basicstyle=\ttfamily\small,
  breaklines=true,
  frame=single,
  framerule=0pt,
  showstringspaces=false,
  columns=fullflexible,
}

% ----- graphics / color -----
\usepackage{graphicx}
\usepackage{xcolor}

% ----- links -----
\usepackage[colorlinks=true,linkcolor=blue!60!black,citecolor=blue!60!black,urlcolor=blue!60!black]{hyperref}
\usepackage{url}

% ----- references -----
\usepackage[round,authoryear]{natbib}

% ----- compact lists -----
\usepackage{enumitem}
\setlist[itemize]{noitemsep,topsep=2pt}
\setlist[enumerate]{noitemsep,topsep=2pt}

\title{Structured Task Decomposition Improves Reliability of \\
       LLM-Based Knowledge Analysis}
\author{NnnHU\\
        Independent Researcher\\
        \texttt{47146850+NnnHU@users.noreply.github.com}}
\date{\today}

\begin{document}
\maketitle

\begin{abstract}
Large language models often fail silently when asked to produce structured
analytical output in a single call: they return empty responses, unparseable
text, or refuse the task. We compare five execution strategies of increasing
complexity on a fixed 13-case corpus of financial documents, holding the model
constant where possible to isolate each strategy's contribution. We find that
(1) \textbf{task decomposition} (an extract~$\rightarrow$~analyze~$\rightarrow$~judge
pipeline) is the dominant driver of reliability, improving the success rate from
31\% (single shot) to 92\%, whereas retry alone adds only +8 points; (2) a
\textbf{multi-model Panel+Judge} architecture, evaluated under blinded dual-LLM
grading with a human anchor, produces higher-quality output than a single-model
pipeline on accuracy, coverage, overall quality, and hallucination rate. We
release the benchmark harness, a corpus manifest with content fingerprints, and
the grading tooling so the results can be reproduced and extended to other
domains.
\end{abstract}

% =====================================================================
\section{Introduction}
% =====================================================================

A common engineering experience with LLMs: you ask for a structured JSON
analysis and the model returns nothing, or a refusal, or prose with no JSON.
For toy tasks this is a nuisance; for a product that must return reliable
structured output, it is fatal.

Two intuitive fixes are often tried: \textbf{retry on failure}, and
\textbf{add more models} (a ``mixture of agents''). Both sound plausible, but
without a controlled experiment it is unclear which one actually buys
reliability---and at what cost.

This report describes a benchmark designed to \textbf{isolate the contribution
of each factor}. We define five execution strategies (\S\ref{sec:setup}) that
form a chain of single-variable changes, run them on a shared corpus, and
measure both \textbf{reliability} (does valid structured output come back?) and
\textbf{quality} (is that output any good?). The design lets us attribute gains
to specific mechanisms rather than to ``using a fancier setup.''

% =====================================================================
\section{Experimental Setup}
\label{sec:setup}
% =====================================================================

\subsection{The five execution strategies}

Each strategy answers the same question against the same source document and
produces the same target structure (a four-field verdict: consensus /
divergence / blindspots / verdict). They differ only in \emph{how} the work is
organized.

\begin{table}[ht!]
\centering
\small
\renewcommand{\arraystretch}{1.15}
\caption{The five execution strategies. ``Cost'' is the number of API calls.
``Isolates'' names the single factor each step is designed to attribute.}
\label{tab:modes}
\begin{tabularx}{\textwidth}{@{}p{2.5cm}Xcc@{}}
\toprule
\textbf{Mode} & \textbf{What it does} & \textbf{Cost} & \textbf{Isolates} \\
\midrule
\textbf{M1} single-shot & One call; doc + question in the prompt & 1 & baseline \\
\textbf{M1R} single + retry & Same as M1, but retries on empty/bad output ($\leq$4 attempts) & 1--4 & contribution of \textbf{retry} alone \\
\textbf{M2} single-model pipeline & extract~$\rightarrow$~analyze~$\rightarrow$~judge, all one model & 3 & contribution of \textbf{task decomposition} beyond retry \\
\textbf{M3} multi-model Panel+Judge & extract, then 2 models analyze in parallel, then a Judge synthesizes & 4 & contribution of \textbf{multiple models} beyond pipeline \\
\textbf{M4} single-model self-critique & answer~$\rightarrow$~critique~$\rightarrow$~revise, one model & 3 & whether \textbf{self-reflection} matches multi-model \\
\bottomrule
\end{tabularx}
\end{table}

The chain is designed so each step changes \textbf{ideally one variable}:

\begin{lstlisting}
M1 -> M1R  : +retry, nothing else           (retry's contribution)
M1R -> M2  : +pipeline structure            (decomposition's contribution)
M2  -> M3  : +a second model in the panel   (multi-model's contribution)
M4  <-> M3 : self-critique vs multi-model   (does self-reflection suffice?)
\end{lstlisting}

A residual confound (documented in \S\ref{sec:discussion}): M3's Judge receives
two analyses while M2's receives one, so ``model count'' and ``Judge input
richness'' are not fully separated. We report this honestly rather than
overclaim.

\subsection{Corpus}

13 cases, all from a single financial-investing article series (the Jeremy
Grantham / GMO bubble narrative), chosen for narrative richness and the
presence of specific verifiable facts (dates, P/E ratios, AUM figures).
Composition:

\begin{itemize}
  \item 1 English summary document ($\sim$3.3k chars)
  \item 7 standard Chinese (Traditional) ASR transcripts ($\sim$9--12k chars each;
        contain speech-to-text noise---typos, misheard names)
  \item 3 large documents ($\sim$31--34k chars)
  \item 1 super-large document ($\sim$53k chars)
  \item 1 multi-document case (3 docs concatenated, $\sim$32k chars)
\end{itemize}

Total: $\sim$256k characters. The ASR noise is intentional---it tests
robustness to messy real-world input. The corpus is narrow in domain (finance);
this is a known limitation (\S\ref{sec:discussion}).

\subsection{Models}

\begin{itemize}
  \item Panel / Judge: \texttt{deepseek-v4-flash} (via \texttt{api.deepseek.com},
        OpenAI-compatible protocol)
  \item Second panel model: \texttt{deepseek-v4-pro}
  \item Temperature: 0.7 for analysis/judge, 0.5 for extract, 0.5 for critique
  \item Max tokens per call: 2048; timeout: 360s
\end{itemize}

\subsection{Metrics}

\paragraph{Reliability} (all 13 cases, all 5 modes): a mode ``succeeds'' on a
case iff it returns all four verdict fields with \textbf{real} content---parse-fallback
placeholders like ``\textit{(could not parse \ldots)}'' do \textbf{not} count.
(An earlier draft of this analysis incorrectly counted placeholders as success;
the bug was found and the numbers below are corrected. See \S\ref{sec:discussion}
on methodology.)

\paragraph{Quality} (7 cases where both M2 and M3 produced real output): graded
by \textbf{two independent LLM judges} (Gemini and KIMI K2.6) in a \textbf{blinded
A/B setup}---neither judge knew which output was M2 or M3. Position bias was
cancelled by per-case A/B reordering. A \textbf{human anchor} graded 3 cases
blindly to validate the LLM judges. Scoring rubric: factual accuracy (1--5),
hallucination (Y/N), key-info coverage (1--5), overall quality (1--5), and a
forced preference (A / B / tied).

% =====================================================================
\section{Results}
% =====================================================================

\subsection{Reliability---decomposition, not retry, drives success}

\begin{table}[h]
\centering
\small
\caption{Reliability across the five modes (all 13 cases).}
\label{tab:reliability}
\begin{tabular}{@{}lcccc@{}}
\toprule
\textbf{Mode} & \textbf{Success rate} & \textbf{Avg latency} & \textbf{Avg output chars} & \textbf{API calls} \\
\midrule
M1 single-shot                   & \textbf{31\%} (4/13)  & 21.4s & 458   & 13 \\
M1R single + retry               & \textbf{38\%} (5/13)  & 45.7s & 626   & 21 \\
M2 single-model pipeline         & \textbf{92\%} (12/13) & 55.3s & 1{,}339 & 39 \\
M3 multi-model Panel+Judge       & \textbf{92\%} (12/13) & 50.4s & 1{,}531 & 52 \\
M4 single-model self-critique    & \textbf{38\%} (5/13)  & 74.2s & 899   & 39 \\
\bottomrule
\end{tabular}
\end{table}

\paragraph{Reading the chain.}

\begin{itemize}
  \item \textbf{M1~$\rightarrow$~M1R: +8 points.} Retry rescued only 1 of 9
        failures. The empty-response problem is not a transient blip that a
        second attempt fixes---the model genuinely struggles with the one-shot
        structured-JSON prompt $\sim$60\% of the time. \textbf{Retry alone is
        insufficient.}
  \item \textbf{M1R~$\rightarrow$~M2: +54 points.} Decomposing the task into
        extract~$\rightarrow$~analyze~$\rightarrow$~judge (with the same single
        model) takes reliability from 38\% to 92\%. This is the dominant effect.
        \textbf{Task decomposition---not retry---is what makes the pipeline
        reliable.}
  \item \textbf{M2~$\rightarrow$~M3: +0 points on reliability} (both 92\%), but
        see quality (\S3.2). Adding a second model does not further improve the
        \emph{rate} of producing valid output once a pipeline is in place.
  \item \textbf{M4 (self-critique): 38\%.} A single model critiquing and revising
        its own output does \textbf{not} approach the pipeline's reliability.
        This is a naive answer$\rightarrow$critique$\rightarrow$revise
        implementation; engineered self-reflection (Reflexion~\citep{shinn2023reflexion},
        Self-Refine~\citep{madaan2023selfrefine},
        Constitutional AI~\citep{bai2022constitutional}) may do better---but the
        result shows self-critique is not automatically a substitute for
        decomposition.
\end{itemize}

\begin{quote}
\textbf{Headline finding:} \emph{Structured task decomposition dramatically
improves the reliability of LLM-based structured analysis---from $\sim$31\% to
$\sim$92\%. The gain comes from the decomposition, not from retrying.}
\end{quote}

\subsection{Quality---multi-model beats single-model pipeline \\
            (blinded, dual-LLM + human-anchored)}

For the 7 cases where both M2 and M3 returned real analysis, blinded grading by
Gemini + KIMI K2.6 (14 votes total):

\begin{table}[h]
\centering
\small
\caption{Quality grading: M2 (single-model pipeline) vs.\ M3 (multi-model
Panel+Judge) on the 7 cases where both produced real output. 14 votes = 2
judges $\times$ 7 cases.}
\label{tab:quality}
\begin{tabular}{@{}lccc@{}}
\toprule
\textbf{Metric} & \textbf{M2 (single-model)} & \textbf{M3 (multi-model)} & \textbf{$\Delta$ (M3$-$M2)} \\
\midrule
Factual accuracy (1--5)        & 3.14 & \textbf{4.14} & \textbf{+1.00} \\
Key-info coverage (1--5)       & 2.93 & \textbf{3.79} & \textbf{+0.86} \\
Overall quality (1--5)         & 2.86 & \textbf{3.71} & \textbf{+0.86} \\
Hallucination (count /14)      & 3    & \textbf{2}    & M3 fewer \\
Preference wins (of 14)        & 6    & \textbf{8}    & M3 wins \\
\bottomrule
\end{tabular}
\end{table}

\textbf{The strongest single signal: both LLM judges agreed on preference in
7/7 cases (100\% inter-rater agreement).} Two architecturally different models,
grading independently and blinded, picked the same winner every time---hard to
produce by chance, indicating the M2-vs-M3 difference is a real signal, not
scoring-scale noise.

\textbf{Human anchor validation: 3/3 agreement.} A domain-knowledgeable human
graded 3 representative cases blindly (deliberately including 2 cases where the
LLMs picked single-model M2---the counter-intuitive choice). The human agreed
with the LLM judges on all 3, including both M2 wins. This confirms the LLM
judges are not blindly pro-multi-model: they pick single-model when it is
actually better, and the human agrees.

\textbf{Reading:} When both modes succeed, multi-model Panel+Judge produces
\textbf{higher-quality} output than a single-model pipeline---more accurate,
better coverage, fewer hallucinations, and preferred more often. The earlier
appearance that ``M2 $\approx$ M3'' (from output-character counts) was a
measurement artifact: by a richness proxy they looked equal, but on real quality
grading multi-model pulls ahead.

\subsection{Notable secondary findings}

\begin{itemize}
  \item \textbf{Parallelism offsets latency.} M3 (4 calls, 2 models in parallel)
        averaged 50.4s vs.\ M2's 55.3s (3 sequential calls). Multi-model fan-out
        does \textbf{not} cost wall-clock time here, though it costs $\sim$33\%
        more tokens.
  \item \textbf{M3 beat M2 on latency in 6/13 cases}, including the
        multi-document case (M2 101s vs.\ M3 37s) where sequential
        extract~$\rightarrow$~analyze~$\rightarrow$~judge on one model is slower
        than parallel panels.
\end{itemize}

% =====================================================================
\section{Reproducibility}
% =====================================================================

The benchmark is fully reproducible from the accompanying reference
implementation, anchored to repository commit \texttt{9301168}.

\begin{lstlisting}[language=bash]
# Run all 5 modes on all 13 cases (~40-70 min, ~150 API calls)
npm run bench -- --full

# Incremental: run only M1R + M4, merging prior M1/M2/M3 traces
npm run bench

# Re-run only M2 + M3 (e.g. after fixing an extract-retry issue)
npm run bench -- --remediate
\end{lstlisting}

\textbf{Corpus:} \texttt{bench-samples/} (13 cases + \texttt{samples.json}
manifest). The corpus is domain-narrow (finance); replacing it with
other-domain documents tests generalization.

\textbf{Configuration:} \texttt{.env} with \texttt{VITE\_VERDEX\_PROVIDER\_*}
and \texttt{VITE\_VERDEX\_PROVIDER2\_*} (two OpenAI-compatible providers; the
environment variable names are historical and retained for reproducibility).

\textbf{Outputs:} \texttt{bench-results/} (not redistributed):
\begin{itemize}
  \item \verb|<runId>-<case>-<query>-trace.json| --- full-fidelity traces (untruncated)
  \item \verb|<runId>-<case>-<query>-report.md| --- per-case comparison tables
  \item \verb|<runId>-SUMMARY.md| --- aggregate + key-question analysis
\end{itemize}

\textbf{Blinded quality grading pack:} \verb|npx tsx scripts/extract-grading.ts|
generates per-case A/B files (\texttt{bench-results/grading-pack/}) for blinded
grading by any LLM or human. The A/B~$\rightarrow$~mode mapping is kept in
\texttt{quality-grading-key.json} (secret until grading is complete).

\paragraph{Note on corpus distribution.} A reference implementation, the grading
harness, and a corpus \textbf{manifest with content fingerprints} accompany this
submission. The 13 source documents are publicly accessible video transcripts;
we provide SHA-256 fingerprints and character counts for verification rather
than redistributing the transcripts, to respect the original creators'
distribution rights. Readers can reproduce the runs by obtaining the source
documents and verifying them against the published fingerprints.

% =====================================================================
\section{Discussion}
\label{sec:discussion}
% =====================================================================

\subsection{What the evidence supports}

\begin{itemize}
  \item \textbf{Strong:} task decomposition improves reliability (31\%~$\rightarrow$~92\%,
        n=13, clean isolation of retry vs.\ decomposition).
  \item \textbf{Preliminary but consistent:} multi-model Panel+Judge improves
        quality over a single-model pipeline (5 metrics, n=7, dual-LLM 7/7
        agreement, human anchor 3/3). We phrase this as \emph{``evidence
        consistently favors''} rather than \emph{``proves''}---n=7 and a single
        domain are not general proof.
\end{itemize}

\subsection{Limitations}

\begin{itemize}
  \item \textbf{Sample size.} n=13 for reliability, n=7 for quality. Directional,
        not statistically definitive.
  \item \textbf{Domain narrowness.} All cases are financial-investing text.
        Whether the decomposition effect holds across domains (code, legal,
        medical) is untested. The corpus is designed to be swappable to test
        this.
  \item \textbf{One residual confound in M2 vs.\ M3.} M3's Judge receives two
        analyses while M2's receives one, so ``model count'' and ``Judge input
        richness'' co-vary. The result should be read as ``multi-model
        Panel+Judge > single-model pipeline,'' not ``more models alone > fewer
        models alone.'' A fully isolated test would feed M2's Judge two copies
        of the same single-model analysis.
  \item \textbf{M2's low scores are partly extract-driven.} In 2--3 cases M2's
        extract step partially failed even after retry, biasing its quality
        scores downward. The direction (M3 > M2) holds on the clean cases too,
        but the gap magnitude may be inflated.
  \item \textbf{LLM judges, not only humans.} The human anchor (3 cases)
        validates the LLM judges' direction, but a larger human study would
        strengthen this.
\end{itemize}

\subsection{Methodological note: the placeholder bug}

An earlier draft of this analysis reported ``100\% success'' for every mode.
That was a statistics bug: the parser emits non-empty placeholder strings
(``\textit{(could not parse \ldots)}'') on failure, and a naive non-empty check
counted them as success. The bug was caught by noticing the data looked
implausibly clean, and the fix (excluding placeholders from the success
criterion) yielded the numbers in \S3. We retain this story in the report
because catching and correcting one's own measurement bias is a core part of
trustworthy benchmarking---and a reminder that \emph{the most dangerous failures
are the ones that look like success}.

\subsection{Implications for practice}

For teams building LLM products that must return reliable structured output:

\begin{enumerate}
  \item \textbf{Decompose first.} Before adding more models or fancier prompts,
        split the task into extract~$\rightarrow$~analyze~$\rightarrow$~judge.
        This single change had the largest effect in our benchmark.
  \item \textbf{Retry is cheap insurance but not a fix.} Add it, but don't expect
        it to solve a fundamentally hard one-shot prompt.
  \item \textbf{Multi-model earns its cost on quality, not reliability.} Once a
        pipeline is reliable, a second model improves output quality (accuracy,
        coverage, hallucination) without adding latency (parallelism offsets it).
  \item \textbf{Self-critique is not a free substitute.} A naive
        answer$\rightarrow$critique$\rightarrow$revise loop did not match the
        pipeline. If pursuing self-reflection, use engineered prompts
        (Reflexion-style~\citep{shinn2023reflexion}), and benchmark it.
\end{enumerate}

% =====================================================================
\section{Related Work}
\label{sec:related}
% =====================================================================

\begin{itemize}
  \item \textbf{Self-Refine / Reflexion / Constitutional AI}---
        single-model self-improvement loops
        \citep{madaan2023selfrefine,shinn2023reflexion,bai2022constitutional}.
        Our M4 is a naive instance; engineered variants may outperform it. The
        benchmark's M4 slot is where such methods can be plugged in for
        comparison.
  \item \textbf{Mixture-of-Agents (MoA)}---multi-model aggregation
        \citep{wang2024mixtureofagents}. Our M3 is a minimal MoA (2 panels + 1
        judge). The benchmark isolates MoA's contribution by comparing to M2
        (single-model pipeline).
  \item \textbf{LLM-as-judge / human-anchored evaluation}---we use dual-LLM
        blinded grading plus a human anchor \citep{zheng2023llmjudge}, following
        emerging best practice for evaluation when full human studies are
        infeasible.
  \item \textbf{Structured output reliability.} Recent benchmarks evaluate
        constrained-decoding and schema-conformance for structured generation
        \citep{geng2025jsonschemabench}; our work complements these by studying
        \emph{task-structure} decomposition (pipeline design) rather than
        decoding-time constraints as the lever for reliable structured output.
  \item \textbf{Stepwise / decomposed reasoning.} Chain-of-thought prompting
        \citep{wei2022chainofthought} elicits reasoning by asking a single model
        to emit intermediate steps \emph{within one call}. Our extract~$\rightarrow$~analyze~$\rightarrow$~judge
        decomposition is complementary but distinct: it splits the task across
        \emph{separate} calls (and, in M3, separate models), each producing a
        structured artifact that the next step consumes. The reliability gains we
        measure come from this explicit inter-step handoff---which converts silent
        single-shot failures into recoverable per-stage failures---rather than from
        intra-prompt chain-of-thought.
\end{itemize}

% =====================================================================
\section*{Acknowledgments}
% =====================================================================

The experimental design, benchmark implementation, data collection, and
analysis were conducted entirely by the author. The prose of this manuscript
was prepared with assistance from an AI writing tool for drafting and
polishing; all technical content, experimental results, and conclusions were
authored, verified, and are the sole responsibility of the author. Any errors
are the author's own.

% =====================================================================
\appendix
\section{Per-case data}
% =====================================================================

Reliability (all 13 cases $\times$ M1/M1R/M2/M3/M4) and quality grading (7 cases
$\times$ M2/M3 $\times$ 2 judges) are available as machine-readable traces in
\verb|bench-results/<runId>-*-trace.json|. The aggregate summary is in
\verb|bench-results/<runId>-SUMMARY.md|. A condensed per-case table:

\begin{table}[h]
\centering
\small
\caption{Per-case reliability. $\checkmark$ = produced real four-field output;
$\times$ = empty / placeholder / refusal. ``Real'' excludes parse-fallback
placeholders.}
\label{tab:percase}
\begin{tabular}{@{}lcccccc@{}}
\toprule
\textbf{Case} & \textbf{Corpus (chars)} & \textbf{M1} & \textbf{M1R} & \textbf{M2} & \textbf{M3} & \textbf{M4} \\
\midrule
grantham-summary     & 3.3k  & $\checkmark$ & $\checkmark$ & $\checkmark$ & $\checkmark$ & $\checkmark$ \\
g1-mean-reversion    & 12.3k & $\checkmark$ & $\checkmark$ & $\checkmark$ & $\checkmark$ & $\times$      \\
g2-keystone          & 10.4k & $\times$     & $\times$     & $\times$     & $\times$     & $\checkmark$ \\
g3-ai-bet            & 9.2k  & $\times$     & $\checkmark$ & $\checkmark$ & $\checkmark$ & $\times$      \\
g4-reinvest          & 11.0k & $\checkmark$ & $\times$     & $\checkmark$ & $\times$     & $\times$      \\
g5-showdown          & 10.1k & $\checkmark$ & $\checkmark$ & $\checkmark$ & $\times$     & $\times$      \\
g6-derivatives       & 10.2k & $\checkmark$ & $\checkmark$ & $\checkmark$ & $\checkmark$ & $\checkmark$ \\
g7-chanos            & 11.1k & $\times$     & $\checkmark$ & $\checkmark$ & $\checkmark$ & $\times$      \\
fbig1-1999           & 31.9k & $\times$     & $\checkmark$ & $\checkmark$ & $\times$     & $\checkmark$ \\
fbig2-2009           & 31.4k & $\times$     & $\times$     & $\times$     & $\times$     & $\times$      \\
fbig3-2026           & 33.8k & $\times$     & $\times$     & $\times$     & $\times$     & $\times$      \\
gsuper-synthesis     & 53.0k & $\checkmark$ & $\checkmark$ & $\checkmark$ & $\times$     & $\times$      \\
multi-g123 ($\times$3) & 31.9k & $\times$   & $\times$     & $\times$     & $\checkmark$ & $\times$      \\
\bottomrule
\end{tabular}
\end{table}

% =====================================================================
\section{Blinded grading protocol}
% =====================================================================

\begin{enumerate}
  \item For each of the 7 cases where both M2 and M3 produced real output,
        extract the four-field verdicts into a case file with A/B labels assigned
        by a stable seeded shuffle (the mapping is saved to
        \texttt{quality-grading-key.json}, kept secret during grading).
  \item Each case file contains: source document, question, Output A, Output B,
        and a scoring rubric.
  \item Feed each case file to two LLM judges (Gemini, KIMI K2.6) independently.
  \item To cancel position bias, the two judges receive different A/B orderings
        per case (the harness varies order by a per-case seed).
  \item Collect accuracy / hallucination / coverage / overall scores + a forced
        preference from each judge.
  \item \textbf{Human anchor:} a human grades 3 cases (selected to include M2-win
        and M3-win cases) using the same blinded files, then the human's
        preferences are compared to the LLM judges' for agreement.
  \item Unblind: map A/B back to M2/M3 via the key file and aggregate.
\end{enumerate}

\bigskip
\noindent\textit{Engineering Report~$\cdot$~reproducible from repository commit
\texttt{9301168}. Corpus, harness, and traces are included. The benchmark is
designed to be re-run on other domains by swapping \texttt{bench-samples/}.}

\bibliographystyle{plainnat}
\bibliography{references}

\end{document}
